% GENERATED FILE. Edit canonical lessons, then run npm run generate:books.
% canonical-source-hash: fnv1a64:ffa4bda300488920
% canonical-lessons: LA-C379-quaestus, LA-C379-stipendium, LA-C379-mercenarius, LA-C379-laboriosus, LA-C379-diligentia

\chapter{Earning a Living}
\label{ch:la-earning-a-living}
\hlchaptermodality{379}

\begin{chapteropening}
\textbf{By the end of this chapter:} \emph{I can talk about earning a living, pay, hired men, hard work and care.}

Complete the last lesson of chapter 379: I can talk about earning a living, pay, hired men, hard work and care.
\end{chapteropening}

\section[quaestus, quaestūs]{quaestus, quaestūs — gain, profit; a way of earning a living}

\label{lesson:LA-C379-quaestus}



\begin{quote}
Before the new one: say the Latin for to sew, then the Latin for to carve.
\end{quote}

\subsection*{You'll want to know: quaestus, quaestūs}
\textbf{quaestus, quaestūs} — \textquotedblleft{}gain, profit; a way of earning a living\textquotedblright{}.

\begin{quote}
Piscātor ex marī quaestum facit.
\end{quote}

\begin{quote}
\emph{The fisherman makes his living from the sea.}
\end{quote}

\emph{quaestus, quaestūs} m., built on \emph{quaerō}, to seek: what you go out to get. \emph{Quaestum facere}, to make money.

\begin{cousinweb}
\emph{quaestus} is the noun built on \emph{quaerō} \textquotedblleft{}to seek, to get\textquotedblright{} (EtymDB; Lewis \& Short): what a man goes out to win.
\end{cousinweb}

\begin{grammarlens}[title={What you've built}]
One more word for telling, explaining and describing at B1.
\end{grammarlens}

\subsection*{Your turn}
\begin{itemize}
\raggedright
  \item \emph{Say it:} \emph{quaestus, quaestūs}
  \item \emph{Say it:} \emph{quaestus, quaestūs}, once more
  \item \emph{Say it:} a sentence of your own with \emph{quaestus, quaestūs}
  \item \emph{Recall:} say the Latin for to sew, then the Latin for to carve, then say \emph{quaestus, quaestūs} again
\end{itemize}

\begin{tcolorbox}[breakable,colback=teal!4,colframe=teal!35!black,title={Before you move on}]
What does quaestus, quaestūs mean? (\textquotedblleft{}Gain, profit; a way of earning a living\textquotedblright{}.) Say it once more.
\end{tcolorbox}
\section[stīpendium, stīpendiī]{stīpendium, stīpendiī — pay, wages (especially a soldier's)}

\label{lesson:LA-C379-stipendium}



\begin{quote}
Before the new one: say the Latin for to carve, then the Latin for gain.
\end{quote}

\subsection*{You'll want to know: stīpendium, stīpendiī}
\textbf{stīpendium, stīpendiī} — \textquotedblleft{}pay, wages (especially a soldier's)\textquotedblright{}.

\begin{quote}
Mīlitēs stīpendium accēpērunt.
\end{quote}

\begin{quote}
\emph{The soldiers received their pay.}
\end{quote}

\emph{stīpendium, stīpendiī} n. A soldier's pay, then a year of service: \emph{vīgintī stīpendia}, twenty campaigns. English \emph{stipend}.

\begin{cousinweb}
\emph{stīpendium} is contracted from \emph{stipipendium}, built from \emph{stips} \textquotedblleft{}a small coin\textquotedblright{} and \emph{pendō} \textquotedblleft{}to weigh out, to pay\textquotedblright{} (Lewis \& Short): coins weighed out. English \emph{stipend} comes from it.
\end{cousinweb}

\begin{grammarlens}[title={What you've built}]
One more word for telling, explaining and describing at B1.
\end{grammarlens}

\subsection*{Your turn}
\begin{itemize}
\raggedright
  \item \emph{Say it:} \emph{stīpendium, stīpendiī}
  \item \emph{Say it:} \emph{stīpendium, stīpendiī}, once more
  \item \emph{Say it:} a sentence of your own with \emph{stīpendium, stīpendiī}
  \item \emph{Recall:} say the Latin for to carve, then the Latin for gain, then say \emph{stīpendium, stīpendiī} again
\end{itemize}

\begin{tcolorbox}[breakable,colback=teal!4,colframe=teal!35!black,title={Before you move on}]
What does stīpendium, stīpendiī mean? (\textquotedblleft{}Pay, wages (especially a soldier's)\textquotedblright{}.) Say it once more.
\end{tcolorbox}
\section[mercēnārius, mercēnāriī]{mercēnārius, mercēnāriī — a hired worker; a mercenary}

\label{lesson:LA-C379-mercenarius}



\begin{quote}
Before the new one: say the Latin for gain, then the Latin for pay.
\end{quote}

\subsection*{You'll want to know: mercēnārius, mercēnāriī}
\textbf{mercēnārius, mercēnāriī} — \textquotedblleft{}a hired worker; a mercenary\textquotedblright{}.

\begin{quote}
Mercēnāriī prō mercēde labōrant.
\end{quote}

\begin{quote}
\emph{Hired men work for a wage.}
\end{quote}

\emph{mercēnārius, mercēnāriī} m. (also spelled \emph{mercennārius}), built on \emph{mercēs}, a wage. As an adjective, hired. English \emph{mercenary}.

\begin{cousinweb}
\emph{mercēnārius} is built on \emph{mercēs, mercēdis} \textquotedblleft{}wages, hire\textquotedblright{} (Lewis \& Short), which belongs with \emph{merx} \textquotedblleft{}goods\textquotedblright{}. English \emph{mercenary} comes from it.
\end{cousinweb}

\begin{grammarlens}[title={What you've built}]
One more word for telling, explaining and describing at B1.
\end{grammarlens}

\subsection*{Your turn}
\begin{itemize}
\raggedright
  \item \emph{Say it:} \emph{mercēnārius, mercēnāriī}
  \item \emph{Say it:} \emph{mercēnārius, mercēnāriī}, once more
  \item \emph{Say it:} a sentence of your own with \emph{mercēnārius, mercēnāriī}
  \item \emph{Recall:} say the Latin for gain, then the Latin for pay, then say \emph{mercēnārius, mercēnāriī} again
\end{itemize}

\begin{tcolorbox}[breakable,colback=teal!4,colframe=teal!35!black,title={Before you move on}]
What does mercēnārius, mercēnāriī mean? (\textquotedblleft{}A hired worker; a mercenary\textquotedblright{}.) Say it once more.
\end{tcolorbox}
\section[labōriōsus, labōriōsa, labōriōsum]{labōriōsus, labōriōsa, labōriōsum — toilsome; hard-working}

\label{lesson:LA-C379-laboriosus}



\begin{quote}
Before the new one: say the Latin for pay, then the Latin for a hired worker.
\end{quote}

\subsection*{You'll want to know: labōriōsus, labōriōsa, labōriōsum}
\textbf{labōriōsus, labōriōsa, labōriōsum} — \textquotedblleft{}toilsome; hard-working\textquotedblright{}.

\begin{quote}
Vīta agricolae labōriōsa est.
\end{quote}

\begin{quote}
\emph{A farmer's life is full of toil.}
\end{quote}

\emph{labōriōsus, labōriōsa, labōriōsum}: of a task, toilsome; of a person, hard-working. English \emph{laborious}.

\begin{cousinweb}
\emph{labōriōsus} is built on \emph{labor, labōris} \textquotedblleft{}toil\textquotedblright{} (EtymDB; Lewis \& Short), with the ending \emph{-ōsus} \textquotedblleft{}full of\textquotedblright{}. English \emph{laborious} comes from it.
\end{cousinweb}

\begin{grammarlens}[title={What you've built}]
One more word for telling, explaining and describing at B1.
\end{grammarlens}

\subsection*{Your turn}
\begin{itemize}
\raggedright
  \item \emph{Say it:} \emph{labōriōsus, labōriōsa, labōriōsum}
  \item \emph{Say it:} \emph{labōriōsus, labōriōsa, labōriōsum}, once more
  \item \emph{Say it:} a sentence of your own with \emph{labōriōsus, labōriōsa, labōriōsum}
  \item \emph{Recall:} say the Latin for pay, then the Latin for a hired worker, then say \emph{labōriōsus, labōriōsa, labōriōsum} again
\end{itemize}

\begin{tcolorbox}[breakable,colback=teal!4,colframe=teal!35!black,title={Before you move on}]
What does labōriōsus, labōriōsa, labōriōsum mean? (\textquotedblleft{}Toilsome; hard-working\textquotedblright{}.) Say it once more.
\end{tcolorbox}
\section[dīligentia, dīligentiae]{dīligentia, dīligentiae — care, diligence}

\label{lesson:LA-C379-diligentia}



\begin{quote}
Before the new one: say the Latin for a hired worker, then the Latin for toilsome.
\end{quote}

\subsection*{You'll want to know: dīligentia, dīligentiae}
\textbf{dīligentia, dīligentiae} — \textquotedblleft{}care, diligence\textquotedblright{}.

\begin{quote}
Opus magnā dīligentiā perfēcit.
\end{quote}

\begin{quote}
\emph{He finished the work with great care.}
\end{quote}

\emph{dīligentia, dīligentiae} f., the noun of \emph{dīligēns}, careful, hard-working. English \emph{diligence}.

\begin{cousinweb}
\emph{dīligentia} is the noun built on \emph{dīligēns} \textquotedblleft{}careful\textquotedblright{} (Lewis \& Short), the participle of a verb made from \emph{dis-} and \emph{legō} \textquotedblleft{}to pick out\textquotedblright{}. English \emph{diligence} comes from it.
\end{cousinweb}

\begin{grammarlens}[title={What you've built}]
One more word for telling, explaining and describing at B1.
\end{grammarlens}

\subsection*{Your turn}
\begin{itemize}
\raggedright
  \item \emph{Say it:} \emph{dīligentia, dīligentiae}
  \item \emph{Say it:} \emph{dīligentia, dīligentiae}, once more
  \item \emph{Say it:} a sentence of your own with \emph{dīligentia, dīligentiae}
  \item \emph{Recall:} say the Latin for a hired worker, then the Latin for toilsome, then say \emph{dīligentia, dīligentiae} again
\end{itemize}

\begin{tcolorbox}[breakable,colback=teal!4,colframe=teal!35!black,title={Before you move on}]
What does dīligentia, dīligentiae mean? (\textquotedblleft{}Care, diligence\textquotedblright{}.) Say it once more.
\end{tcolorbox}
